\section{Introduzione}
Nel presente capitolo viene descritto in dettaglio il piano di lavoro preventivamente delineato per lo svolgimento dello stage presso il Dipartimento di Matematica ``Tullio Levi-Civita'' dell'Università degli Studi di Padova. Tale attività formativa si propone di approfondire tematiche di ricerca operativa applicate alla gestione e ottimizzazione dello spazio aereo europeo, attraverso l'utilizzo e l'analisi di dati forniti da \citet{eurocontrol2025} tramite il repository \acrfull{adrr}.

\section{Finalità dello Stage}
L'obiettivo principale dello stage consiste nell'analisi e nell'elaborazione di dati riguardanti il traffico e la configurazione dello spazio aereo europeo. Tali dati presentano informazioni storiche e diversificate circa le configurazioni adottate e la domanda di traffico aereo nei cieli europei. Le attività svolte durante lo stage sono orientate alla comprensione approfondita di tali dati e allo sviluppo di procedure di esportazione e analisi per facilitare la loro applicazione in modelli di ottimizzazione per la configurazione dinamica dello spazio aereo.

\newpage
\section{Ripartizione ore}
\subsection{Preventivo}
La pianificazione delle attività, definita nel preventivo iniziale, prevedeva la distribuzione delle ore di lavoro riportata nella Tabella~\ref{tab:pianificazione-attivita}.
\begin{center}
    \input{chapters/ch2/tabella_ripartizione_ore}
\end{center}
\clearpage

\section{Obbiettivi}
\label{sec:obbiettivi}
\subsection{Notazione}
Si farà riferimento ai requisiti secondo le seguenti notazioni:
\begin{itemize}
	\item \textit{O} per i requisiti obbligatori, vincolanti in quanto obiettivo primario richiesto dal committente
	\item \textit{D} per i requisiti desiderabili, non vincolanti o strettamente necessari,
		  ma dal riconoscibile valore aggiunto
	\item \textit{F} per i requisiti facoltativi, rappresentanti valore aggiunto non strettamente 
		  competitivo
\end{itemize}

Le sigle precedentemente indicate saranno seguite da una coppia sequenziale di numeri, identificativo del requisito.

\DisablePunct 
\subsection{Obiettivi fissati}
Si prevede lo svolgimento dei seguenti obiettivi:
\begin{itemize}
	\item Obbligatori
	\begin{itemize}
		\item \underline{\textit{O01}}: Acquisizione di competenze su dati concernenti il traffico aereo e il suo controllo;
	    \item \underline{\textit{O02}}: Acquisizione dei dati \acrshort{adrr};
	    \item \underline{\textit{O03}}: Implementazione di procedure di esportazione dei dati \acrshort{adrr} nei formati richiesti dai modelli di ottimizzazione per \acrshort{dac};
	    \item \underline{\textit{O04}}: Redazione della documentazione di un manuale utente per le procedure realizzate.
	\end{itemize}
	
	\item Desiderabili 
	\begin{itemize}
	   \item \underline{\textit{D01}}: Realizzazione di un benchmark di istanze reali DAC;
	   \item \underline{\textit{D02}}: Test di performance dei modelli di ottimizzazione per \acrshort{dac} sul benchmark realizzato.	
	\end{itemize}
	\item Facoltativi
	\begin{itemize}
	   \item \underline{\textit{F01}}: Implementazione di un'interfaccia utente per l'esecuzione delle procedure di importazione ed esportazione dei dati \acrshort{adrr}.
	\end{itemize} 
\end{itemize}
\EnablePunct


%TODO da scrivere di più
\section{Risorse utilizzate}
Durante lo svolgimento delle attività dello stage, si è fatto uso delle seguenti risorse strumentali e software:

\begin{itemize}
\item Ambiente Python per l'elaborazione e analisi dei dati
\item Librerie specializzate per ottimizzazione (Cplex)
\item Repository EUROCONTROL per i dati storici e configurazioni
\item Strumenti per la gestione della documentazione (\LaTeX e Git)
\end{itemize}

\newpage